\section{Empirical Overview Across NC, WI, TX, and FL}
\label{sec:empirical}

\subsection{Data and Method}

All empirical results use VEST 2020 presidential precinct returns mapped to
2020 census tract geometry via areal interpolation. Presidential returns are
used for consistency across states and years; the 2020 presidential election
provides the most recent high-turnout baseline before the 2022 congressional
maps took effect. Enacted map results refer to the 2022 congressional plans
in effect for the 2022 elections in each state.

\textsc{Bisect} maps are generated using the following configuration:
\begin{itemize}
\item Structure: \texttt{standard-bisect}
\item Weights: \texttt{--weights-override geographic}
\item Search: single seed (results marked with $^\dagger$)
\item Census year: 2020
\end{itemize}

Partisan metrics are computed from district-level aggregated vote shares using
the uniform swing model for EG, Bias, and the seats-votes curve. MM and
Declination are computed directly from observed district vote shares without
the swing assumption.

\subsection{Six-Metric Comparison Table}

Table~\ref{tab:six-metric} presents all six metrics for \textsc{bisect}
algorithmic maps and enacted 2022 congressional maps across the four states.
All \textsc{bisect} values are single-run results ($^\dagger$).

\begin{table}[ht]
\centering
\caption{Six-Metric Partisan Fairness Comparison: Bisect vs.\ Enacted (2022)}
\label{tab:six-metric}
\begin{threeparttable}
\small
\begin{tabular}{@{}llrrrrrr@{}}
\toprule
State & Plan & EG & MM & Bias & Decl.\ $\delta$ & Resp.\ $R$ & LSq \\
\midrule
NC ($k=14$) & Bisect$^\dagger$ & $+0.02$ & $+0.02$ & $-0.01$ & $+0.08$ & $2.0$ & $0.05$ \\
NC ($k=14$) & Enacted 2022    & $+0.09$ & $+0.07$ & $-0.07$ & $+0.35$ & $1.3$ & $0.14$ \\
\addlinespace
WI ($k=8$)  & Bisect$^\dagger$ & $+0.02$ & $+0.02$ & $-0.01$ & $+0.07$ & $1.9$ & $0.04$ \\
WI ($k=8$)  & Enacted 2022    & $+0.06$ & $+0.06$ & $-0.05$ & $+0.28$ & $1.4$ & $0.11$ \\
\addlinespace
TX ($k=38$) & Bisect$^\dagger$ & $+0.03$ & $+0.03$ & $-0.02$ & $+0.09$ & $2.1$ & $0.07$ \\
TX ($k=38$) & Enacted 2022    & $+0.11$ & $+0.08$ & $-0.09$ & $+0.40$ & $1.2$ & $0.18$ \\
\addlinespace
FL ($k=28$) & Bisect$^\dagger$ & $+0.02$ & $+0.02$ & $-0.01$ & $+0.08$ & $2.0$ & $0.06$ \\
FL ($k=28$) & Enacted 2022    & $+0.08$ & $+0.06$ & $-0.06$ & $+0.32$ & $1.3$ & $0.13$ \\
\bottomrule
\end{tabular}
\begin{tablenotes}
\small
\item $^\dagger$ Single-run \textsc{bisect} result. EG$>0$ indicates Republican
advantage; Bias$<0$ indicates Republican advantage; Decl.\ $\delta>0$ indicates
Republican advantage through cracking; Resp.\ $R$: higher is more competitive;
LSq: lower is more proportional.
\item Source: VEST 2020 presidential precinct returns; 2020 Census tract geometry;
enacted maps as codified for 2022 elections.
\end{tablenotes}
\end{threeparttable}
\end{table}

\subsection{Cross-State Patterns}

Several patterns are consistent across all four states:

\textbf{Magnitude reduction.} \textsc{Bisect} maps reduce the absolute value
of all four descriptive metrics by at least 50\% relative to enacted maps in
every state. The reduction is largest for Declination (bisect $|\delta| < 0.10$
vs.\ enacted $|\delta| \approx 0.30$--$0.40$) because Declination is most
sensitive to the geometric packing-cracking configurations that compactness
directly prevents.

\textbf{Sign consistency.} All metrics have the same sign direction (favorable
to Republicans) for both \textsc{bisect} and enacted maps in all four states,
reflecting the underlying geographic sorting of Democratic voters into dense
urban cores. The key finding is not that \textsc{bisect} maps are perfectly
neutral---they are not---but that the residual partisan skew is attributable
to geography rather than deliberate manipulation.

\textbf{Responsiveness improvement.} \textsc{Bisect} maps achieve $R \approx
2.0$ (competitive ideal) compared to enacted maps at $R \approx 1.2$--$1.4$.
This improvement reflects the creation of genuinely competitive districts
that are insulated from incumbent protection in enacted maps.

\textbf{Proportionality improvement.} Gallagher LSq improves by roughly 60\%
across all four states. The residual LSq in \textsc{bisect} maps ($0.04$--$0.07$)
represents the structural floor imposed by single-member district geography;
the reduction from enacted values ($0.11$--$0.18$) represents the mapmaker
contribution that \textsc{bisect} eliminates.

\subsection{The Gap as Evidence}

The gap between \textsc{bisect} and enacted values on each metric provides
the empirical foundation for state court claims. In \textit{Harper v.\ Hall},
the North Carolina Supreme Court held that the state constitution's ``free
elections'' clause (Art.\ I, \S 10) prohibits maps that ``prioritize'' partisan
advantage ``to the exclusion of all else.'' The six-metric comparison in
Table~\ref{tab:six-metric} translates this standard into measurable terms:
if a neutral algorithm produces NC maps with EG $= 0.02$ and the enacted map
has EG $= 0.09$, the gap of $0.07$ is not explained by geography (which affects
both equally) but by deliberate mapmaker choices.

The strength of this argument does not depend on the court adopting any
particular metric or threshold. It depends only on the court accepting that
(1) the \textsc{bisect} maps are genuinely neutral (demonstrable by the
algorithm's lack of partisan data inputs), and (2) the enacted maps are
further from the neutral centroid than natural geographic sorting explains
(demonstrated by the six-metric table). The convergence of all six metrics
in the same direction, and the consistency of the finding across four states,
makes metric-shopping criticism untenable: the result is not an artifact
of metric selection.
